\section{Manutencao e Operacoes}

\subsection{Rotinas de Manutencao}

\subsubsection{Manutencao Diaria}

\begin{lstlisting}[caption=Script de manutencao diaria]
#!/bin/bash
# manutencao_diaria.sh

DB_NAME="finops"
LOGFILE="/var/log/finops_daily.log"

echo "$(date): Iniciando manutencao diaria" >> $LOGFILE

# 1. Processar recorrencias vencidas
echo "$(date): Processando recorrencias" >> $LOGFILE
RECORRENCIAS_PROCESSADAS=$(psql -d $DB_NAME -t -c "SELECT processar_recorrencias();")
echo "$(date): $RECORRENCIAS_PROCESSADAS recorrencias processadas" >> $LOGFILE

# 2. Atualizar saldos das contas
echo "$(date): Atualizando saldos" >> $LOGFILE
psql -d $DB_NAME -c "
DO \$\$
DECLARE
    usuario_record RECORD;
BEGIN
    FOR usuario_record IN SELECT id FROM usuarios WHERE ativo = true
    LOOP
        CALL atualizar_saldos_usuario(usuario_record.id);
    END LOOP;
END \$\$;
"

# 3. Atualizar metas de poupanca
echo "$(date): Atualizando progresso das metas" >> $LOGFILE
psql -d $DB_NAME -c "
UPDATE metas SET 
    valor_atual = (
        SELECT COALESCE(SUM(t.valor), 0)
        FROM transacoes t
        WHERE t.categoria_id = metas.categoria_id
          AND t.usuario_id = metas.usuario_id
          AND t.tipo = 'receita'
          AND t.status = 'confirmada'
    ),
    alcancada = (valor_atual >= valor_alvo),
    data_alcancada = CASE 
        WHEN valor_atual >= valor_alvo AND data_alcancada IS NULL 
        THEN CURRENT_TIMESTAMP 
        ELSE data_alcancada 
    END
WHERE ativa = true;
"

# 4. Limpar sessoes expiradas
echo "$(date): Limpando sessoes expiradas" >> $LOGFILE
SESSOES_REMOVIDAS=$(psql -d $DB_NAME -t -c "
    DELETE FROM sessoes 
    WHERE data_expiracao < CURRENT_TIMESTAMP 
       OR data_ultimo_acesso < CURRENT_TIMESTAMP - INTERVAL '7 days';
    SELECT ROW_COUNT();
")
echo "$(date): $SESSOES_REMOVIDAS sessoes removidas" >> $LOGFILE

# 5. Refresh das views materializadas criticas
echo "$(date): Atualizando views materializadas" >> $LOGFILE
psql -d $DB_NAME -c "REFRESH MATERIALIZED VIEW CONCURRENTLY mv_resumo_mensal;"

echo "$(date): Manutencao diaria concluida" >> $LOGFILE
\end{lstlisting}

\subsubsection{Manutencao Mensal}

\begin{lstlisting}[caption=Script de manutencao mensal]
#!/bin/bash
# manutencao_mensal.sh

DB_NAME="finops"
LOGFILE="/var/log/finops_monthly.log"

echo "$(date): Iniciando manutencao mensal" >> $LOGFILE

# 1. Arquivar dados antigos
echo "$(date): Arquivando dados antigos" >> $LOGFILE

# Criar tabela de arquivo se nao existir
psql -d $DB_NAME -c "
CREATE TABLE IF NOT EXISTS arquivo.transacoes_antigas (
    LIKE transacoes INCLUDING ALL
);
"

# Mover transacoes com mais de 2 anos para arquivo
psql -d $DB_NAME -c "
INSERT INTO arquivo.transacoes_antigas 
SELECT * FROM transacoes 
WHERE data_transacao < CURRENT_DATE - INTERVAL '2 years';

DELETE FROM transacoes 
WHERE data_transacao < CURRENT_DATE - INTERVAL '2 years';
"

# 2. Recriar indices fragmentados
echo "$(date): Recriando indices fragmentados" >> $LOGFILE
psql -d $DB_NAME -c "
DO \$\$
DECLARE
    idx_record RECORD;
BEGIN
    FOR idx_record IN 
        SELECT schemaname, tablename, indexname
        FROM pg_stat_user_indexes
        WHERE idx_scan > 1000 AND idx_tup_read > idx_tup_fetch * 2
    LOOP
        EXECUTE 'REINDEX INDEX ' || quote_ident(idx_record.indexname);
    END LOOP;
END \$\$;
"

# 3. Atualizar estatisticas completas
echo "$(date): Atualizando estatisticas completas" >> $LOGFILE
psql -d $DB_NAME -c "ANALYZE VERBOSE;"

# 4. Verificar integridade referencial
echo "$(date): Verificando integridade referencial" >> $LOGFILE
psql -d $DB_NAME -c "
-- Verificar orfaos em transacoes
SELECT 'Transacoes sem usuario:' as problema, count(*) as quantidade
FROM transacoes t
LEFT JOIN usuarios u ON t.usuario_id = u.id
WHERE u.id IS NULL

UNION ALL

SELECT 'Transacoes sem conta:' as problema, count(*) as quantidade  
FROM transacoes t
LEFT JOIN contas c ON t.conta_id = c.id
WHERE c.id IS NULL

UNION ALL

SELECT 'Contas sem usuario:' as problema, count(*) as quantidade
FROM contas c
LEFT JOIN usuarios u ON c.usuario_id = u.id
WHERE u.id IS NULL AND c.compartilhada = false;
"

# 5. Otimizar tamanho do banco
echo "$(date): Executando VACUUM FULL em tabelas grandes" >> $LOGFILE
psql -d $DB_NAME -c "
SELECT 'VACUUM FULL ' || schemaname || '.' || tablename || ';'
FROM pg_tables 
WHERE schemaname = 'public'
  AND pg_total_relation_size(schemaname||'.'||tablename) > 100*1024*1024; -- >100MB
" | psql -d $DB_NAME

echo "$(date): Manutencao mensal concluida" >> $LOGFILE
\end{lstlisting}

\subsection{Monitoramento Operacional}

\subsubsection{Dashboard de Status}

\begin{lstlisting}[caption=Funcao para status geral do sistema]
CREATE OR REPLACE FUNCTION status_sistema()
RETURNS TABLE (
    componente TEXT,
    status TEXT,
    detalhes TEXT,
    ultimo_check TIMESTAMP
) AS $$
BEGIN
    -- Status do banco de dados
    RETURN QUERY SELECT 
        'Banco de Dados'::TEXT,
        'OK'::TEXT,
        'Conectado e funcionando'::TEXT,
        CURRENT_TIMESTAMP;
    
    -- Status das conexoes
    RETURN QUERY 
    WITH conexoes AS (
        SELECT count(*) as total FROM pg_stat_activity WHERE state = 'active'
    )
    SELECT 
        'Conexoes Ativas'::TEXT,
        CASE WHEN total > 80 THEN 'ALERTA' ELSE 'OK' END::TEXT,
        'Total: ' || total::TEXT,
        CURRENT_TIMESTAMP
    FROM conexoes;
    
    -- Status do tamanho do banco
    RETURN QUERY
    WITH tamanho AS (
        SELECT pg_size_pretty(pg_database_size('finops')) as size_text,
               pg_database_size('finops') as size_bytes
    )
    SELECT 
        'Tamanho BD'::TEXT,
        CASE WHEN size_bytes > 10*1024*1024*1024 THEN 'ATENCAO' ELSE 'OK' END::TEXT,
        size_text,
        CURRENT_TIMESTAMP
    FROM tamanho;
    
    -- Status dos backups
    RETURN QUERY
    WITH ultimo_backup AS (
        SELECT 
            EXTRACT(EPOCH FROM (CURRENT_TIMESTAMP - MAX(timestamp_operacao)))/3600 as horas
        FROM audit.audit_logs 
        WHERE tabela = 'sistema' AND operacao = 'BACKUP'
    )
    SELECT 
        'Ultimo Backup'::TEXT,
        CASE 
            WHEN horas IS NULL THEN 'ERRO'
            WHEN horas > 48 THEN 'ATRASADO' 
            ELSE 'OK' 
        END::TEXT,
        CASE 
            WHEN horas IS NULL THEN 'Nenhum backup encontrado'
            ELSE 'Ha ' || round(horas::numeric, 1) || ' horas'
        END,
        CURRENT_TIMESTAMP
    FROM ultimo_backup;
    
    -- Status das recorrencias
    RETURN QUERY
    WITH recorrencias_pendentes AS (
        SELECT count(*) as pendentes
        FROM recorrencias 
        WHERE ativa = true AND proxima_execucao < CURRENT_DATE
    )
    SELECT 
        'Recorrencias'::TEXT,
        CASE WHEN pendentes > 10 THEN 'ATENCAO' ELSE 'OK' END::TEXT,
        pendentes::TEXT || ' pendentes',
        CURRENT_TIMESTAMP
    FROM recorrencias_pendentes;
END;
$$ LANGUAGE plpgsql;
\end{lstlisting}

\subsubsection{Alertas Automaticos}

\begin{lstlisting}[caption=Sistema de alertas baseado em triggers]
-- Tabela para armazenar alertas
CREATE TABLE sistema.alertas (
    id BIGSERIAL PRIMARY KEY,
    tipo VARCHAR(50) NOT NULL,
    severidade INTEGER NOT NULL, -- 1=info, 2=warning, 3=error, 4=critical
    titulo VARCHAR(200) NOT NULL,
    descricao TEXT,
    usuario_id UUID REFERENCES usuarios(id),
    resolvido BOOLEAN DEFAULT false,
    data_criacao TIMESTAMP DEFAULT CURRENT_TIMESTAMP,
    data_resolucao TIMESTAMP
);

-- Funcao para criar alerta
CREATE OR REPLACE FUNCTION criar_alerta(
    p_tipo TEXT,
    p_severidade INTEGER,
    p_titulo TEXT,
    p_descricao TEXT DEFAULT NULL,
    p_usuario_id UUID DEFAULT NULL
)
RETURNS void AS $$
BEGIN
    INSERT INTO sistema.alertas (tipo, severidade, titulo, descricao, usuario_id)
    VALUES (p_tipo, p_severidade, p_titulo, p_descricao, p_usuario_id);
    
    -- Se for critico, notificar imediatamente
    IF p_severidade >= 4 THEN
        PERFORM pg_notify('alerta_critico', 
            json_build_object(
                'tipo', p_tipo,
                'titulo', p_titulo,
                'usuario_id', p_usuario_id
            )::text
        );
    END IF;
END;
$$ LANGUAGE plpgsql;

-- Trigger para detectar transacoes suspeitas
CREATE OR REPLACE FUNCTION verificar_transacao_suspeita()
RETURNS TRIGGER AS $$
DECLARE
    media_usuario DECIMAL(15,2);
    limite_alerta DECIMAL(15,2);
BEGIN
    -- Calcular media de transacoes do usuario nos ultimos 30 dias
    SELECT COALESCE(AVG(valor), 0) INTO media_usuario
    FROM transacoes 
    WHERE usuario_id = NEW.usuario_id 
      AND tipo = NEW.tipo
      AND data_transacao >= CURRENT_DATE - INTERVAL '30 days';
    
    limite_alerta := media_usuario * 5; -- 5x a media
    
    -- Criar alerta se valor for muito alto
    IF NEW.valor > limite_alerta AND limite_alerta > 100 THEN
        PERFORM criar_alerta(
            'transacao_alta',
            2,
            'Transacao com valor elevado detectada',
            'Valor: R$ ' || NEW.valor || ' (Media: R$ ' || round(media_usuario, 2) || ')',
            NEW.usuario_id
        );
    END IF;
    
    RETURN NEW;
END;
$$ LANGUAGE plpgsql;

CREATE TRIGGER trg_verificar_transacao_suspeita
    AFTER INSERT ON transacoes
    FOR EACH ROW EXECUTE FUNCTION verificar_transacao_suspeita();
\end{lstlisting}

\subsection{Procedures de Operacao}

\subsubsection{Gerenciamento de Usuarios}

\begin{lstlisting}[caption=Procedures para operacoes de usuario]
-- Procedure para desativar usuario
CREATE OR REPLACE PROCEDURE desativar_usuario(p_usuario_id UUID)
LANGUAGE plpgsql AS $$
BEGIN
    -- Desativar usuario
    UPDATE usuarios SET ativo = false WHERE id = p_usuario_id;
    
    -- Desativar contas
    UPDATE contas SET ativa = false WHERE usuario_id = p_usuario_id;
    
    -- Encerrar sessoes ativas
    UPDATE sessoes SET ativa = false WHERE usuario_id = p_usuario_id;
    
    -- Log da operacao
    INSERT INTO audit.audit_logs (tabela, operacao, usuario_id, dados_novos)
    VALUES ('usuarios', 'DEACTIVATE', p_usuario_id, 
            jsonb_build_object('timestamp', CURRENT_TIMESTAMP));
    
    COMMIT;
END;
$$;

-- Procedure para migrar dados entre usuarios
CREATE OR REPLACE PROCEDURE migrar_dados_usuario(
    p_usuario_origem UUID,
    p_usuario_destino UUID
)
LANGUAGE plpgsql AS $$
BEGIN
    -- Verificar se usuarios existem
    IF NOT EXISTS (SELECT 1 FROM usuarios WHERE id = p_usuario_origem) THEN
        RAISE EXCEPTION 'Usuario origem nao encontrado';
    END IF;
    
    IF NOT EXISTS (SELECT 1 FROM usuarios WHERE id = p_usuario_destino) THEN
        RAISE EXCEPTION 'Usuario destino nao encontrado';
    END IF;
    
    -- Migrar transacoes
    UPDATE transacoes SET usuario_id = p_usuario_destino 
    WHERE usuario_id = p_usuario_origem;
    
    -- Migrar contas (marcar como transferidas)
    UPDATE contas SET 
        usuario_id = p_usuario_destino,
        nome = nome || ' (Transferida)'
    WHERE usuario_id = p_usuario_origem;
    
    -- Migrar categorias personalizadas
    UPDATE categorias SET usuario_id = p_usuario_destino 
    WHERE usuario_id = p_usuario_origem AND sistema = false;
    
    -- Log da migracao
    INSERT INTO audit.audit_logs (tabela, operacao, dados_novos)
    VALUES ('usuarios', 'MIGRATE', 
            jsonb_build_object(
                'origem', p_usuario_origem,
                'destino', p_usuario_destino,
                'timestamp', CURRENT_TIMESTAMP
            ));
    
    COMMIT;
END;
$$;
\end{lstlisting}

\subsection{Automacao e Jobs}

\subsubsection{Configuracao de Jobs}

\begin{lstlisting}[caption=Configuracao de cron jobs]
# /etc/cron.d/finops

# Manutencao diaria (2:00 AM)
0 2 * * * postgres /opt/finops/scripts/manutencao_diaria.sh

# Backup diario (3:00 AM)
0 3 * * * postgres /opt/finops/scripts/backup_completo.sh

# Verificacao de backup (3:30 AM)
30 3 * * * postgres /opt/finops/scripts/verificar_backups.sh

# Manutencao semanal (domingo 4:00 AM)
0 4 * * 0 postgres /opt/finops/scripts/manutencao_semanal.sh

# Manutencao mensal (primeiro domingo 5:00 AM)
0 5 1-7 * 0 postgres /opt/finops/scripts/manutencao_mensal.sh

# Monitoramento de performance (a cada 15 minutos)
*/15 * * * * postgres /opt/finops/scripts/alertas_performance.sh

# Processamento de recorrencias (todo dia 1:00 AM)
0 1 * * * postgres psql -d finops -c "SELECT processar_recorrencias();"

# Limpeza de logs (semanal, sabado 23:00)
0 23 * * 6 postgres /opt/finops/scripts/limpeza_logs.sh
\end{lstlisting}
